\Chapter{49}

\Verse{1}
{
    \zA{话说香菱见众人正说笑他，便迎上去笑道：“你们看这首诗：要使得，我就还 学； 要还不好，我就死了这做诗的心了。” 说着，把诗递与黛玉及众人看时，只见
        写道是： 精华欲掩料应难，影自娟娟魄自寒。 一片砧敲千里白，半轮鸡唱五更残。 绿蓑江上秋闻笛，红袖楼头夜倚栏。 博得嫦娥应自问：何缘不使永团？？
        众人看了，笑道：“这首不但好，而且新巧有意趣。 可知俗语说：‘天下无难事， 只怕有心人。’ 社里一定请你了！”}
}
{
    \eA{Hsiang Ling, we will now proceed, perceived the young ladies engaged in
        chatting and laughing, and went up to them with a smiling countenance. "Just
        you look at this stanza!" she said. "If it's all right, then I'll continue
        my studies; but if it isn't worth any thing, I'll banish at once from my
        mind all idea of going in for versification." With these words, she handed
        the verses to Tai-yü and her companions.}
}
{
}

\Verse{2}
{
    \zA{香菱听了，心下不信，料着是他们哄自己的话， 还只管问黛玉宝钗等。 正说之间，只见几个小丫头并老婆子忙忙的走来，都笑道：“来了好些姑娘奶
        奶们，我们都不认得； 奶奶姑娘们快认亲去。” 李纨笑道：“这是那里的话？ 你到 底说明白了，是谁的亲戚？” 那婆子丫头都笑道：“奶奶的两位妹子都来了；
        还有 一位姑娘，说是薛大姑娘的妹子； 还有一位爷，说是薛大爷的兄弟。 我这会子请姨 太太去呢，奶奶和姑娘们先上去罢。” 说着，一径去了。}
}
{
    \eA{When they came to look at them, they found this to be their burden:  If thou
        would'st screen Selene's beauteous sheen, thou'lt find it  hard. Her shadows
        are by nature full of grace, frigid her form. A row of clothes-stones
        batter, while she lights a thousand li. When her disc's half, and the cock
        crows at the fifth watch, 'tis  cold. Wrapped in my green cloak in autumn, I
        hear flutes on the stream.}
}
{
}

\Verse{3}
{
    \zA{宝钗笑道：“我们薛蝌和 他妹子来了不成？” 李纨笑道：“或者我婶娘又上京来了？ 怎么他们都凑在一处？ 这 可是奇事。”
        大家来至王夫人上房，只见黑压压的一地。 又有邢夫人的嫂子，带了女儿岫烟 进京来投邢夫人的，可巧凤姐之兄王仁也正进京，两亲家一处搭帮来了。 走至半路
        泊船时，遇见李纨寡婶，带着两个女儿，长名李纹，次名李绮，也上京，大家叙起 来，又是亲戚，因此三家一路同行。}
}
{
    \eA{While in the tower the red-sleeved maid leans on the rails at night. She
        feels also constrained to ask of the goddess Ch'ang O:  'Why it is that she
        does not let the moon e'er remain round?' "This stanza is not only good,"
        they with one voice exclaimed, after perusing it, "but it's original, it's
        charming. It bears out the proverb: 'In the world, there's nothing
        difficult;}
}
{
}

\Verse{4}
{
    \zA{后有薛蟠之从弟薛蝌，因当年父亲在京时，已 将胞妹薛宝琴许配都中梅翰林之子为妻，正欲进京聘嫁，闻得王仁进京，他也随后 带了妹子赶来。
        所以今日会齐了，来访投各人亲戚。 于是大家见礼叙过，贾母王夫 人都欢喜非常。 贾母因笑道：“怪道昨日晚上灯花爆了又爆，结了又结，原来应到 今日。”
        一面叙些家常，收了带来的礼物，一面命留酒饭。 凤姐儿自不必说，忙上 加忙； 李纨宝钗自然和婶母姊妹叙离别之情。}
}
{
    \eA{the only thing hard to get at is a human being with a will.' We'll certainly
        ask you to join our club." Hsiang Ling caught this remark; but so little did
        she credit it that fancying that they were making fun of her, she still went
        on to press Tai-yü, Pao-ch'ai and the other girls to give her their
        opinions.}
}
{
}

\Verse{5}
{
    \zA{黛玉见了，先是欢喜，后想起众人皆 有亲眷，独自己孤单无倚，不免又去垂泪。 宝玉深知其情，十分劝慰了一番方罢。
        然后宝玉忙忙来至怡红院中，向袭人、麝月、晴雯笑道：“你们还不快着看去! 谁知宝姐姐的亲哥哥是那个样子，他这叔伯兄弟，形容举止另是个样子，倒像是宝
        姐姐的同胞兄弟似的。 更奇在你们成日家只说宝姐姐是绝色的人物，你们如今瞧见 他这妹子，还有大嫂子的两个妹子，我竟形容不出来了。}
}
{
    \eA{But while engaged in speaking, she spied a number of young waiting-maids,
        and old matrons come with hurried step. "Several young ladies and ladies
        have come," they announced smilingly, "but we don't know any of them. So
        your ladyship and you, young ladies, had better come at once and see what
        relatives they are." "What are you driving at?" Li Wan laughed.}
}
{
}

\Verse{6}
{
    \zA{老天，老天，你有多少精 华灵秀，生出这些人上之人来!可知我‘井底之蛙’，成日家只说现在的这几个人
        是有一无二的，谁知不必远寻，就是本地风光，一个赛似一个。 如今我又长了一层 学问了。 除了这几个，难道还有几个不成？” 一面说，一面自笑。
        袭人见他又有些 魔意，便不肯去瞧。 晴雯等早去瞧了一遍回来，带笑向袭人说道：“你快瞧瞧去!
        大太太一个侄女儿，宝姑娘一个妹妹，大奶奶两个妹妹，倒像一把子四根水葱儿。”}
}
{
    \eA{"You might, after all, state distinctly whose relatives they are." "Your
        ladyship's two young sisters have come," the matrons and maids rejoined
        smiling. "There's also another young lady, who says she's miss Hsüeh's
        cousin, and a gentleman who pretends to be Mr. Hsüeh P'an's junior cousin.
        We are now off to ask Mrs. Hsüeh to meet them. But your ladyship and the
        young ladies might go in advance and greet them."}
}
{
}

\Verse{7}
{
    \zA{一语未了，只见探春也笑着进来找宝玉，因说：“咱们诗社可兴旺了。” 宝玉 笑道：“正是呢。 这是一高兴起诗社，鬼使神差来了这些人。 但只一件，不知他们
        可学过做诗不曾？” 探春道：“我才都问了问，虽是他们自谦，看其光景，没有不 会的。 便是不会也没难处，你看香菱就知道了。” 晴雯笑道：“他们里头薛大姑娘
        的妹妹更好。 三姑娘看着怎么样？” 探春道：“果然的。 据我看来，连他姐姐并这 些人总不及他。”}
}
{
    \eA{As they spoke, they straightway took their leave. "Has our Hsüeh K'o come
        along with his sisters?" Pao-ch'ai inquired, with a smile. "My aunt has
        probably also come to the capital," Li Wan laughed. "How is it they've all
        arrived together? This is indeed a strange thing!" Then adjourning in a body
        into Madame Wang's drawing rooms, they saw the floor covered with a black
        mass of people.}
}
{
}

\Verse{8}
{
    \zA{袭人听了，又是诧异，又笑道：“这也奇了，还从那里再寻好的 去呢？ 我倒要瞧瞧去。” 探春道：“老太太一见了，喜欢的无可不可的，已经逼着
        咱们太太认了干女孩儿了。 老太太要养活，才刚已经定了。” 宝玉喜的忙问：“这 话果然么？” 探春道：“我几时撒过谎？”
        又笑道：“老太太有了这个好孙女儿， 就忘了你这孙子了。” 宝玉笑道：“这倒不妨，原该多疼女孩儿些是正理。 明儿十 六，咱们可该起社了。”}
}
{
    \eA{Madame Hsing's sister-in-law was there as well. She had entered the capital
        with her daughter, Chou Yen, to look up madame Hsing. But lady Feng's
        brother, Wang Jen, had, as luck would have it, just been preparing to start
        for the capital, so the two family connexions set out in company for their
        common destination.}
}
{
}

\Verse{9}
{
    \zA{探春道：“林丫头刚起来了，二姐姐又病了，终是七上八 下的。” 宝玉道：“二姐姐又不大做诗，没有他又何妨。” 探春道：“索性等几天，
        等他们新来的混熟了，咱们邀上他们岂不好？ 这会子大嫂子宝姐姐心里自然没有诗 兴的。 况且湘云没来，颦儿才好了，人都不合式。 不如等着云丫头来了，这几个新
        的也熟了，颦儿也大好了，大嫂子和宝姐姐心也闲了，香菱诗也长进了：如此邀一 满社，岂不好？}
}
{
    \eA{After accomplishing half their journey, they encountered, while their boats
        were lying at anchor, Li Wan's widowed sister-in-law, who also was on her
        way to the metropolis, with her two girls, the elder of whom was Li Wen and
        the younger Li Ch'i. They all them talked matters over, and, induced by the
        ties of relationship, the three families prosecuted their voyage together.}
}
{
}

\Verse{10}
{
    \zA{咱们两个如今且往老太太那里去听听，除宝姐姐的妹妹不算外，他 一定是在咱们家住定了的。 倘或那三个要不在咱们这里住，咱们央告着老太太，留
        下他们也在园子里住了，咱们岂不多添几个人，越发有趣了。” 宝玉听了，喜的眉开眼笑，忙说道：“倒是你明白。 我终久是个糊涂心肠，空
        喜欢了一会子，却想不到这上头。” 说着，兄妹两个一齐往贾母处来。 果然王夫人 已认了薛宝琴做干女儿，贾母喜欢非常，不命往园中住，晚上跟着贾母一处安寝。}
}
{
    \eA{But subsequently, Hsüeh P'an's cousin Hsüeh K'o,—whose father had, when on a
        visit years ago to the capital, engaged his uterine sister to the son of the
        Han-lin Mei, whose residence was in the metropolis,—came while planning to
        go and consummate the marriage, to learn of Wang Jen's departure, so taking
        his sister with him, he kept in his track till he managed to catch him up.}
}
{
}

\Verse{11}
{
    \zA{薛蝌自向薛蟠书房中住下了。 贾母和邢夫人说：“你侄女儿也不必家去了，园里住 几天，逛逛再去。” 邢夫人兄嫂家中原艰难，这一上京原仗的是邢夫人与他们治房
        舍、帮盘缠，听如此说，岂不愿意。 邢夫人便将邢岫烟交与凤姐儿。 凤姐儿算着园 中姊妹多，性情不一，且又不便另设一处，莫若送到迎春一处去，倘日后邢岫烟有
        些不遂意的事，纵然邢夫人知道了，与自己无干。}
}
{
    \eA{Hence it happened that they all now arrived in a body to look up their
        respective relatives. In due course, they exchanged the conventional
        salutations; and these over, they had a chat. Dowager lady Chia and madame
        Wang were both filled with ineffable delight. "Little wonder is it," smiled
        old lady Chia, "if the snuff of the lamp crackled time and again; and if it
        formed and reformed into a head!}
}
{
}

\Verse{12}
{
    \zA{从此后，若邢岫烟家去住的日期 不算，若在大观园住到一个月上，凤姐儿亦照迎春分例，送一分与岫烟。 凤姐儿冷 眼？？
        岫烟心性行为，竟不像邢夫人及他的父母一样，却是个极温厚可疼的人。 因 此凤姐儿反怜他家贫命苦，比别的姊妹多疼他些，邢夫人倒不大理论了。 贾母王夫
        人等因素喜李纨贤惠，且年轻守节，令人敬服，今见他寡婶来了，便不肯叫他外头 去住。 那婶母虽十分不肯，无奈贾母执意不从，只得带着李纹李绮在稻香村住下了。}
}
{
    \eA{It was, indeed, sure to come to this to-day!" While she conversed on
        every-day topics, the presents had to be put away; and, as she, at the same
        time, expressed a wish to keep the new arrivals to partake of some wine and
        eatables, lady Feng had, needless to say, much extra work added to her
        ordinary duties.}
}
{
}

\Verse{13}
{
    \zA{当下安插既定，谁知忠靖侯史鼎又迁委了外省大员，不日要带家眷去上任，贾 母因舍不得湘云，便留下他了，接到家中。 原要命凤姐儿另设一处与他住，史湘云
        执意不肯，只要和宝钗一处住，因此也就罢了。 此时大观园中，比先又热闹了多少：李纨为首，馀者迎春、探春、惜春、宝钗、
        黛玉、湘云、李纹、李绮、宝琴、邢岫烟，再添上凤姐儿和宝玉，一共十三人。 叙 起年庚，除李纨年纪最长，凤姐次之，馀者皆不过十五六七岁，大半同年异月，连
        他们自己也不能记清谁长谁幼；}
}
{
    \eA{Li Wan and Pao-ch'ai descanted, of course, with their aunts and cousins on
        the events that had transpired since their separation.}
}
{
}

\Verse{14}
{
    \zA{并贾母王夫人及家中婆子丫头也不能细细分清，不 过是“姐”“妹”“兄”“弟”四个字，随便乱叫。
        如今香菱正满心满意只想做诗，又不敢十分罗唆宝钗，可巧来了个史湘云，那 史湘云极爱说话的，那里禁得香菱又请教他谈诗？ 越发高了兴，没昼没夜，高谈阔
        论起来。 宝钗因笑道：“我实在聒噪的受不得了。 一个女孩儿家，只管拿着诗做正 经事讲起来，叫有学问的人听了反笑话，说不守本分。}
}
{
    \eA{But Tai-yü, though when they first met, continued in cheerful spirits, could
        not again, when the recollection afterwards flashed through her mind that
        one and all had their relatives, and that she alone had not a soul to rely
        upon, avoid withdrawing out of the way, and giving vent to tears.}
}
{
}

\Verse{15}
{
    \zA{一个香菱没闹清，又添上你 这个话口袋子，满口里说的是什么：怎么是‘杜工部之沉郁，韦苏州之淡雅’，又 怎么是‘温八叉之绮靡，李义山之隐僻’。
        痴痴癫癫，那里还像两个女儿家呢？” 说得香菱湘云二人都笑起来。 正说着，只见宝琴来了，披着一领斗篷，金翠辉煌，不知何物。 宝钗忙问：“这 是那里的？”
        宝琴笑道：“因下雪珠儿，老太太找了这一件给我的。” 香菱上来瞧 道：“怪道这么好看，原来是孔雀毛织的。” 湘云笑道：“那里是孔雀毛？}
}
{
    \eA{Pao-yü, however, read her feelings, and he had to do all that lay in his
        power to exhort her and to console her for a time before she cheered up.
        Pao-yü then hurried into the I Hung court. Going up to Hsi Jen, She Yüeh and
        Chi'ng Wen: "Don't you yet hasten to go and see them?" he smiled. "Who'd
        ever have fancied that cousin Pao-ch'ai's own cousin would be what he is?}
}
{
}

\Verse{16}
{
    \zA{就是野 鸭子头上的毛做的。 可见老太太疼你了：这么着疼宝玉，也没给他穿。” 宝钗笑道： “真是俗语说的，‘各人有各人的缘法’。
        我也想不到他这会子来，既来了，又有 老太太这么疼他。” 湘云道：“你除了在老太太跟前，就在园里，来这两处，只管 玩笑吃喝。
        到了太太屋里，若太太在屋里，只管和太太说笑，多坐一回无妨； 若太 太不在屋里，你别进去。 那屋里人多心坏，都是耍咱们的。” 说的宝钗、宝琴、香
        菱、莺儿等都笑了。}
}
{
    \eA{That cousin of hers is so unique in appearance and in deportment. He looks
        as if he were cousin Pao-ch'ai's uterine younger brother. But what's still
        more odd is, that you should have kept on saying the whole day long that
        cousin Pao-ch'ai is a very beautiful creature. You should now see her
        cousin, as well as the two girls of her senior sister-in-law. I couldn't
        adequately tell you what they're like. Good heavens!}
}
{
}

\Verse{17}
{
    \zA{宝钗笑道：“说你没心却有心，虽然有心，到底嘴太直了。 我 们这琴儿，今儿你竟认他做亲妹妹罢。” 湘云又瞅了宝琴笑道：“这一件衣裳也只
        配他穿，别人穿了实在不配。” 正说着，只见琥珀走来，笑道：“老太太说了：叫 宝姑娘别管紧了琴姑娘，他还小呢，让他爱怎么着就由他怎么着，他要什么东西只
        管要，别多心。” 宝钗忙起身答应了，又推宝琴笑道：“你也不知是那里来的这点 福气!你倒去罢，恐怕我们委屈了你!我就不信，我那些儿不如你？”}
}
{
    \eA{Good heavens! What subtle splendour and spiritual beauty must you possess to
        produce beings like them, so superior to other human creatures! How plain it
        is that I'm like a frog wallowing at the bottom of a well! I've throughout
        every hour of the day said to myself that nowhere could any girls be found
        to equal those at present in our home; but, as it happens, I haven't had far
        to look!}
}
{
}

\Verse{18}
{
    \zA{说话之间，宝玉黛玉进来了，宝钗犹自嘲笑。 湘云因笑道：“宝姐姐，你这话 虽是玩，却有人真心是这样想呢。” 琥珀笑道：“真心恼的再没别人，就只是他。”
        口里说，手指着宝玉。 宝钗湘云都笑道：“他倒不是这样人。” 琥珀又笑道：“不 是他，就是他。” 说着，又指黛玉。 湘云便不作声。 宝钗笑道：“更不是了。
        我的 妹妹和他的妹妹一样，他喜欢的比我还甚呢，他那里还恼？ 你信云儿混说，他那嘴 有什么正经。”}
}
{
    \eA{Even in our own native sphere, one would appear to eclipse the other! Here I
        have now managed to add one more stratum to my store of learning! But can it
        possibly be that outside these few, there can be any more like them?" As he
        uttered these sentiments, he smiled to himself.}
}
{
}

\Verse{19}
{
    \zA{宝玉素昔深知黛玉有些小性儿，尚不知近日黛玉和宝钗之事，正恐 贾母疼宝琴，他心中不自在。 今儿湘云如此说了，宝钗又如此答，再审度黛玉声色
        亦不似往日，果然与宝钗之说相符，心中甚是不解。 因想：“他两个素日不是这样 的，如今看来，竟更比他人好了十倍。” 一时又见林黛玉赶着宝琴叫“妹妹”，并
        不提名道姓，真似亲姊妹一般。 那宝琴年轻心热，且本性聪敏，自幼读书识字，今 在贾府住了两日，大概人物已知；}
}
{
    \eA{But Hsi Jen noticed how much under the influence of his insane fits he once
        more was, and she promptly abandoned all idea of going over to pay her
        respects to the visitors. Ch'ing Wen and the other girls had already gone
        and seen them and come back. Putting on a smile, "You'd better," they urged
        Hsi Jen, "be off at once and have a look at them.}
}
{
}

\Verse{20}
{
    \zA{又见众姊妹都不是那轻薄脂粉，且又和姐姐皆和 气，故也不肯怠慢。 其中又见林黛玉是个出类拔萃的，便更与黛玉亲敬异常。 宝玉 看着，只是暗暗的纳罕。
        一时宝钗姊妹往薛姨妈房内去后，湘云往贾母处来，林黛玉回房歇着。 宝玉便 找了黛玉来，笑道：“我虽看了《西厢记》，也曾有明白的几句说了取笑，你还曾 恼过。
        如今想来，竟有一句不解，我念出来，你讲讲我听。” 黛玉听了，便知有文 章，因笑道：“你念出来我听听。”}
}
{
    \eA{Our elder mistress' niece, Miss Pao's cousin, and our senior lady's two
        sisters resemble a bunch of four leeks so pretty are they!" But scarcely
        were these words out of their lips, than they perceived T'an Ch'un too enter
        the room, beaming with smiles. She came in quest of Pao-yü. "Our poetical
        society is in a flourishing way," she remarked. "It is," smiled Pao-yü.}
}
{
}

\Verse{21}
{
    \zA{宝玉笑道：“那《闹简》上有一句说的最好： ‘是几时孟光接了梁鸿案？’ 这五个字不过是现成的典，难为他‘是几时’三个虚 字，问的有趣。 是几时接了？
        你说说我听听。” 黛玉听了，禁不住也笑起来，因笑 道：“这原问的好。 他也问的好，你也问的好。” 宝玉道：“先时你只疑我，如今 你也没的说了。”
        黛玉笑道：“谁知他竟真是个好人，我素日只当他藏奸。” 因把 说错了酒令，宝钗怎样说他，连送燕窝，病中所谈之事，细细的告诉宝玉，宝玉方 知原故。}
}
{
    \eA{"Here no sooner do we, in the exuberance of our spirits, start a poetical
        society, than the devils and gods bring through their agency, all these
        people in our midst! There's only one thing however. Have they, I wonder,
        ever learnt how to write poetry or not?" "I just now asked every one of
        them," T'an Ch'un replied.}
}
{
}

\Verse{22}
{
    \zA{因笑道：“我说呢!正纳闷‘是几时孟光接了梁鸿案’，原来是从‘小孩 儿家口没遮拦’上就接了案了。” 黛玉因又说起宝琴来，想起自己没有姊妹，不免又哭了。
        宝玉忙劝道：“这又 自寻烦恼了。 你瞧瞧，今年比旧年越发瘦了，你还不保养。 每天好好的，你必是自 寻烦恼，哭一会子，才算完了这一天的事。”
        黛玉拭泪道：“近来我只觉心酸，眼 泪却像比旧年少了些的。 心里只管酸痛，眼泪却不多。” 宝玉道：“这是你哭惯了， 心里疑惑，岂有眼泪会少的！”}
}
{
    \eA{"Their ideas of themselves are modest, it's true, yet from all I can gather
        there's not one who can't versify. But should there even be any who can't,
        there's nothing hard about it. Just look at Hsiang Ling. Her case will show
        you the truth of what I say." "Of the whole lot," smiled Ch'ing Wen, "Miss
        Hsüeh's cousin carries the palm. What do you think about her, Miss Tertia?"
        "It's really so!" T'an Ch'un responded.}
}
{
}

\Verse{23}
{
    \zA{正说着，只见他屋里的小丫头子送了猩猩毡斗篷来，又说：“大奶奶才打发人 来说：下了雪，要商议明日请人做诗呢。” 一语未了，只见李纨的丫头走来请黛玉。
        宝玉便邀着黛玉同往稻香村来。 黛玉换上掐金挖云红香羊皮小靴，罩了一件大红羽 绉面白狐狸皮的鹤氅，系一条青金闪绿双环四合如意绦，上罩了雪帽。 二人一齐踏
        雪行来，只见众姊妹都在那里，都是一色大红猩猩毡与羽毛缎斗篷，独李纨穿一件 哆罗呢对襟褂子，薛宝钗穿一件莲青斗纹锦上添花洋线番？ 丝的鹤氅。}
}
{
    \eA{"In my own estimation, even her elder cousin and all this bevy of girls are
        not fit to hold a candle to her!" Hsi Jen felt much surprise at what she
        heard. "This is indeed odd!" she smiled. "Whence could one hunt up any
        better? We'd like to go and have a peep at her." "Our venerable senior,"
        T'an Ch'un observed, "was at the very first sight of her so charmed with her
        that there's nothing she wouldn't do.}
}
{
}

\Verse{24}
{
    \zA{邢岫烟仍是 家常旧衣，并没避雨之衣。 一时湘云来了，穿着贾母给他的一件貂鼠脑袋面子、大
        毛黑灰鼠里子、里外发烧大褂子，头上带着一顶挖云鹅黄片金里子大红猩猩毡昭君 套，又围着大貂鼠风领。 黛玉先笑道：“你们瞧瞧，孙行者来了。 他一般的拿着雪
        褂子，故意妆出个小骚鞑子样儿来。” 湘云笑道：“你们瞧我里头打扮的。” 一面 说，一面脱了褂子，只见他里头穿着一件半新的靠色三厢领袖秋香色盘金五色绣龙
        窄？}
}
{
    \eA{She has already compelled our Madame Hsing to adopt her as a godchild. Our
        dear ancestor wishes to bring her up herself; this point was settled a
        little while back." Pao-yü went into ecstasies. "Is this a fact?" he eagerly
        inquired. "How often have I gone in for yarns?" T'an Ch'un said.}
}
{
}

\Verse{25}
{
    \zA{小袖掩衿银鼠短袄，里面短短的一件水红妆缎狐肷褶子，腰里紧紧束着一条蝴 蝶结子长穗五色宫绦，脚下也穿着鹿皮小靴，越显得蜂腰猿背，鹤势螂形。 众人笑
        道：“偏他只爱打扮成个小子的样儿，原比他打扮女儿更俏丽了些。” 湘云笑道：“快商议做诗。 我听听是谁的东家？” 李纨道：“我的主意。 想来
        昨儿的正日已自过了，再等正日还早呢，可巧又下雪，不如咱们大家凑个热闹，又 给他们接风，又可以做诗。 你们意思怎么样？”}
}
{
    \eA{"Now that our worthy senior," continuing, she laughed, "has got this nice
        granddaughter, she has banished from her mind all thought of a grandson like
        you!" "Never mind," answered Pao-yü smiling. "It's only right that girls
        should be more doated upon. But to-morrow is the sixteenth, so we should
        have a meeting." "That girl Lin Tai-yü is no sooner out of bed," T'an Ch'un
        remarked, "than cousin Secunda falls ill again.}
}
{
}

\Verse{26}
{
    \zA{宝玉先道：“这话很是，只是今儿 晚了，若到明儿，晴了又无趣。” 众人都道：“这雪未必晴。 纵晴了，这一夜下的 也够赏了。”
        李纨道：“我这里虽然好，又不如芦雪庭好。 我已经打发人笼地炕去 了，咱们大家拥炉做诗。 老太太想来未必高兴。 况且咱们小玩意儿，单给凤丫头个
        信儿就是了。 你们每人一两银子就够了，送到我这里来。”}
}
{
    \eA{Everything is, in fact, up and down!" "Our cousin Secunda," Pao-yü
        explained, "doesn't also go in very much for verses, so, what would it
        matter if she were left out?" "It would be well to wait a few days," T'an
        Ch'un urged, "until the new comers have had time to see enough of us to
        become intimate. We can then invite them to join us. Won't this be better?}
}
{
}

\Verse{27}
{
    \zA{指着香菱、宝琴、李纹、 李绮、岫烟，“五个不算外，咱们里头二丫头病了不算，四丫头告了假也不算，你 们四分子送了来，我包管五六两银子也尽够了。”
        宝钗等一齐应诺。 因又拟题限韵， 李纨笑道：“我心里早已定了。 等到了明日临期，横竖知道。” 说毕，大家又说了 一回闲话，方往贾母处来，当日无话。
        到了次日清早，宝玉因心里惦记着，这一夜没好生得睡，天亮了就爬起来。}
}
{
    \eA{Our senior sister-in-law and cousin Pao have now no mind for poetry.
        Besides, Hsiang-yün has not arrived. P'in Erh is just over her sickness. The
        members are not all therefore in a fit state, so wouldn't it be preferable
        if we waited until that girl Yün came? The new arrivals will also have a
        chance of becoming friendly. P'in Erh will likewise recover entirely.}
}
{
}

\Verse{28}
{
    \zA{掀 起帐子一看，虽然门窗尚掩，只是窗上光辉夺目，心内早踌躇起来，埋怨定是晴了， 日光已出。
        一面忙起来揭起窗屉，从玻璃窗内往外一看，原来不是日光，竟是一夜 的雪，下的将有一尺厚，天上仍是搓绵扯絮一般。 宝玉此时喜欢非常，忙唤起人来，
        盥漱已毕，只穿一件茄色哆罗呢狐狸皮袄，罩一件海龙小鹰膀褂子，束了腰，披上 玉针蓑，带了金藤笠，登上沙棠屐，忙忙的往芦雪庭来。 出了院门，四顾一望，并
        无二色，远远的是青松翠竹，自己却似装在玻璃盆内一般。}
}
{
    \eA{Our senior sister-in-law and cousin Pao will have time to compose their
        minds; and Hsiang Ling to improve in her verses. We shall then be able to
        convene a full meeting; and won't it be better? You and I must now go over
        to our worthy ancestor's, on the other side, and hear what's up.}
}
{
}

\Verse{29}
{
    \zA{于是走至山坡之下。 顺 着山脚刚转过去，已闻得一股寒香扑鼻，回头一看，却是妙玉那边栊翠庵中有十数 枝红梅如胭脂一般，映着雪色，分外显得精神，好不有趣。
        宝玉便立住，细细的赏 玩了一回方走。 只见蜂腰板桥上一个人打着伞走来，是李纨打发了请凤姐儿去的 人。 宝玉来至芦雪庭，只见丫头婆子正在那里扫雪开径。
        原来这芦雪庭盖在一个傍 山临水河滩之上，一带几间茅檐土壁，横篱竹牖，推窗便可垂钓，四面皆是芦苇掩 覆。}
}
{
    \eA{But, barring cousin Pao-ch'ai's cousin,—for we needn't take her into
        account, as it's sure to have been decided that she should live in our
        home,—if the other three are not to stay here with us, we should entreat our
        grandmother to let them as well take up their quarters in the garden. And if
        we succeed in adding a few more to our number, won't it be more fun for us?"}
}
{
}

\Verse{30}
{
    \zA{一条去径，逶迤穿芦度苇过去，便是藕香榭的竹桥了。 众丫头婆子见他披蓑带 笠而来，都笑道：“我们才说正少一个渔翁，如今果然全了。 姑娘们吃了饭才来呢，
        你也太性急了。” 宝玉听了，只得回来。 刚至沁芳亭，见探春正从秋爽斋出来，围 着大红猩猩毡的斗篷，带着观音兜，扶着个小丫头，后面一个妇人打着一把青绸油
        伞。 宝玉知道他往贾母处去，遂站在亭边等他来到，二人一同出园前去。 宝琴正在里间房内梳洗更衣。 一时众姐妹来齐，宝玉只嚷饿了，连连催饭。}
}
{
    \eA{Pao-yü at these words was so much the more gratified that his very eyebrows
        distended, and his eyes laughed. "You've got your wits about you!" he
        speedily exclaimed. "My mind is ever so dull! I've vainly given way to a fit
        of joy. But to think of these contingencies was beyond me!" So saying the
        two cousins repaired together to their grandmother's suite of apartments;}
}
{
}

\Verse{31}
{
    \zA{好 容易等摆上饭来，头一样菜是牛乳蒸羊羔，贾母就说：“这是我们有年纪人的药， 没见天日的东西，可惜你们小孩子吃不得。
        今儿另外有新鲜鹿肉，你们等着吃罢。” 众人答应了。 宝玉却等不得，只拿茶泡了一碗饭，就着野鸡爪子忙忙的爬拉完了。
        贾母道：“我知道你们今儿又有事情，连饭也不顾吃了。” 就叫：“留着鹿肉给他 晚上吃罢。” 凤姐儿忙说：“还有呢，吃残了的倒罢了。”}
}
{
    \eA{where, in point of fact, Madame Wang had already gone through the ceremony
        of recognizing Hsüeh Pao-ch'in as her godchild. Dowager lady Chia's
        fascination for her, however, was so much out of the common run that she did
        not tell her to take up her quarters in the garden. Of a night, she
        therefore slept with old lady Chia in the same rooms; while Hsüeh K'o put up
        in Hsüeh P'an's study.}
}
{
}

\Verse{32}
{
    \zA{湘云就和宝玉计较道： “有新鹿肉，不如咱们要一块，自己拿了园里弄着，又吃又玩。” 宝玉听了，真和 凤姐要了一块，命婆子送进园去。
        一时大家散后，进园齐往芦雪庭来，听李纨出题限韵。 独不见湘云宝玉二人。 黛玉道：“他两个人再到不得一处，要到了一处，生出多少事来。 这会子一定算计
        那块鹿肉去了。”}
}
{
    \eA{"Your niece needn't either return home," dowager lady Chia observed to
        Madame Hsing. "Let her spend a few days in the garden and see the place
        before she goes." Madame Hsing's brother and sister-in-law were, indeed, in
        straitened circumstances at home.}
}
{
}

\Verse{33}
{
    \zA{正说着，只见李婶娘也走来看热闹，因问李纨道：“怎么那一个 带玉的哥儿和那一个挂金麒麟的姐儿，那样干净清秀，又不少吃的，他两个在那里
        商议着要吃生肉呢，说的有来有去的。 我只不信，肉也生吃得的？” 众人听了，都 笑道：“了不得，快拿了他两个来。” 黛玉笑道：“这可是云丫头闹的。
        我的卦再 不错。” 李纨即忙出来，找着他两个，说道：“你们两个要吃生的，我送你们到老 太太那里吃去，那怕一只生鹿，撑病了不与我相干。}
}
{
    \eA{So much so that they had, on their present visit to the capital, actually to
        rely upon such accommodation as Madame Hsing could procure for them and upon
        such help towards their travelling expenses as she could afford to give
        them. When she consequently heard her proposal, Madame Hsing was, of course,
        only too glad to comply with her wishes, and readily she handed Hsing
        Chou-yen to the charge of lady Feng.}
}
{
}

\Verse{34}
{
    \zA{这么大雪，怪冷的，快替我做 诗去罢。” 宝玉忙笑道：“没有的事!我们烧着吃呢。” 李纨道：“这还罢了。”
        只见老婆子们拿了铁炉、铁叉、铁丝蒙来，李纨道：“留神，割了手不许哭。” 说 着，方进去了。 那边凤姐打发平儿回复不来，为发放年例正忙着呢。
        湘云见了平儿，那里肯放？ 平儿也是个好玩的，素日跟着凤姐儿无所不至，见如此有趣，乐得玩笑，因而退去 手上的镯子，三个人围着火，平儿便要先烧三块吃。}
}
{
    \eA{But lady Feng, bethinking herself of the number of young ladies already in
        the garden, of their divergent dispositions and, above all things, of the
        inconvenience of starting a separate household, deemed it advisable to send
        her to live along with Ying Ch'un;}
}
{
}

\Verse{35}
{
    \zA{那边宝钗黛玉平素看惯了，不 以为异，宝琴等及李婶娘深为罕事。 探春和李纨等已议定了题韵。 探春笑道：“你 们闻闻，香气这里都闻见了，我也吃去。”
        说着，也找了他们来。 李纨也随来，说： “客已齐了，你们还吃不够吗？” 湘云一面吃，一面说道：“我吃这个方爱吃酒， 吃了酒才有诗。
        若不是这鹿肉，今儿断不能做诗。” 说着，只见宝琴披着凫靥裘， 站在那里笑。 湘云笑道：“傻子!你来尝尝。” 宝琴笑道：“怪腌？ 的。”}
}
{
    \eA{for in the event, (she thought), of Hsing Chou-yen meeting afterwards with
        any contrarieties, she herself would be clear of all responsibility, even
        though Madame Hsing came to hear about them.}
}
{
}

\Verse{36}
{
    \zA{宝钗笑 道：“你尝尝去，好吃的很呢。 你林姐姐弱，吃了不消化，不然，他也爱吃。” 宝 琴听了，就过去吃了一块，果然好吃，就也吃起来。
        一时凤姐儿打发小丫头来叫平 儿，平儿说：“史姑娘拉着我呢，你先去罢。” 小丫头去了。 一时，只见凤姐儿也
        披了斗篷走来，笑道：“吃这样好东西，也不告诉我！” 说着，也凑在一处吃起来。 黛玉笑道：“那里找这一群花子去!罢了罢了，今日芦雪庭遭劫，生生被云丫头作
        践了。 我为芦雪庭一大哭。”}
}
{
    \eA{Deducting, therefore any period, spent by Hsing Chou-yen on a visit home,
        lady Feng allowed Hsing Chou-yen as well, if she extended her stay in the
        garden of Broad Vista for any time over a month, an amount equal to that
        allotted to Ying Ch'un. Lady Feng weighed with unprejudiced eye Hsing
        Chou-yen's temperament and deportment.}
}
{
}

\Verse{37}
{
    \zA{湘云冷笑道：“你知道什么!‘是真名士自风流’。 你们都是假清高，最可厌的。 我们这会子腥的膻的大吃大嚼，回来却是锦心绣口。”
        宝钗笑道：“你回来若做的不好了，把那肉掏出来，就把这雪压的芦苇子？ 上些， 以完此劫。” 说着，吃毕，洗了一回手。
        平儿带镯子时，却少了一个，左右前后乱找了一番， 踪迹全无。 众人都诧异。 凤姐儿笑道：“我知道这镯子的去向，你们只管做诗去。
        我们也不用找，只管前头去，不出三日包管就有了。”}
}
{
    \eA{She found in her not the least resemblance to Madame Hsing, or even to her
        father and mother; but thought her a most genial and love-inspiring girl.
        This consideration actuated lady Feng (not to deal harshly with her), but to
        pity her instead for the poverty, in which they were placed at home, and for
        the hard lot she had to bear, and to treat her with far more regard than she
        did any of the other young ladies.}
}
{
}

\Verse{38}
{
    \zA{说着又问：“你们今儿做什 么诗？ 老太太说了，离年又近了，正月里还该做些灯谜儿大家玩笑。” 众人听了， 都笑道：“可是呢，倒忘了。
        如今赶着做几个好的，预备着正月里玩。” 说着，一 齐来至地炕屋内，只见杯盘果菜俱已摆齐了，墙上已贴出诗题、韵脚、格式来了。
        宝玉湘云二人忙看时，只见题目是《即景联句》，“五言排律一首，限‘二萧’韵。” 后面尚未列次序。
        李纨道：“我不大会做诗，我只起三句罢，然后谁先得了谁先联。” 宝钗道：“到底分个次序。”}
}
{
    \eA{Madame Hsing, however, did not lavish much attention on her. Dowager lady
        Chia, Madame Wang and the rest had all along been fond of Li Wan for her
        virtuous and benevolent character. Besides, her continence in remaining a
        widow at her tender age commanded general esteem. When they therefore now
        saw her husbandless sister-in-law come to pay her a visit, they would not
        allow her to go and live outside the mansion.}
}
{
}

\Verse{39}
{
    \zA{要知端底，且看下回分解。 (本章完)}
}
{
    \eA{Her sister-in-law was, it is true, extremely opposed to the proposal, but as
        dowager lady Chia was firm in her determination, she had no other course but
        to settle down, along with Li Wen and Li Ch'i, in the Tao Hsiang village.}
}
{
}

\Verse{40}
{
    \zA{}
}
{
    \eA{They had by this time assigned quarters to all the new comers, when, who
        would have thought it, Shih Ting, Marquis of Chung Ching, was once again
        appointed to a high office in another province, and he had shortly to take
        his family and proceed to his post. But so little could old lady Chia brook
        the separation from Hsiang-yün that she kept her behind and received her in
        her own home. Her original idea was to have asked lady Feng to have separate
        rooms arranged for her, but Shih Hsiang-yün was so obstinate in her refusal,
        her sole wish being to put up with Pao-ch'ai, that the idea had, in
        consequence, to be abandoned. At this period, the garden of Broad Vista was
        again much more full of life than it had ever been before. Li Wan was the
        chief inmate.}
}
{
}

\Verse{41}
{
    \zA{}
}
{
    \eA{The rest consisted of Ying Ch'un, T'an Ch'un, Hsi Ch'un, Pao-ch'ai, Tai-yü,
        Hsiang-yün, Li Wen, Li Ch'i, Pao Ch'in and Hsing Chou-yen. In addition to
        these, there were lady Feng and Pao-yü, so that they mustered thirteen in
        all. As regards age, irrespective of Li Wan, who was by far the eldest, and
        lady Feng, who came next, the other inmates did not exceed fourteen, sixteen
        or seventeen. But the majority of them had come into the world in the same
        year, though in different months, so they themselves could not remember
        distinctly who was senior, and who junior.}
}
{
}

\Verse{42}
{
    \zA{}
}
{
    \eA{Even dowager lady Chia, Madame Wang and the matrons and maids in the
        household were unable to tell the differences between them with any
        accuracy, given as they were to the simple observance of addressing
        themselves promiscuously and quite at random by the four words representing
        'female cousin' and 'male cousin.' Hsiang Ling was gratifying her wishes to
        her heart's content and devoting her mind exclusively to the composition of
        verses, not presuming however to make herself too much of a nuisance to
        Pao-ch'ai, when, by a lucky coincidence, Shih Hsiang-yün came on the scene.
        But how was it possible for one so loquacious as Hsiang-yün to avoid the
        subject of verses, when Hsiang Ling repeatedly begged her for explanations?}
}
{
}

\Verse{43}
{
    \zA{}
}
{
    \eA{This inspirited her so much the more, that not a day went by, yea not a
        single night, on which she did not start some loud argument and lengthy
        discussion. "You really," Pao-ch'ai felt impelled to laugh, "kick up such a
        din, that it's quite unbearable! Fancy a girl doing nothing else than
        turning poetry into a legitimate thing for raising an argument! Why, were
        some literary persons to hear you, they would, instead of praising you, have
        a laugh at your expense, and say that you don't mind your own business. We
        hadn't yet got rid of Hsiang Ling with all her rubbish, and here we have a
        chatterbox like you thrown on us! But what is it that that mouth of yours
        keeps on jabbering? What about the bathos of Tu Kung-pu; and the unadorned
        refinement of Wei Su-chou?}
}
{
}

\Verse{44}
{
    \zA{}
}
{
    \eA{What also about Wen Pa-ch'a's elegant diction; and Li I-shan's abstruseness?
        A pack of silly fools that you are! Do you in any way behave like girls
        should?" These sneers evoked laughter from both Hsiang Ling and Hsiang-yün.
        But in the course of their conversation, they perceived Pao-ch'in drop in,
        with a waterproof wrapper thrown over her, so dazzling with its gold and
        purplish colours, that they were at a loss to make out what sort of article
        it could be. "Where did you get this?" Pao-ch'ai eagerly inquired. "It was
        snowing," Pao-ch'in smilingly replied, "so her venerable ladyship turned up
        this piece of clothing and gave it to me."}
}
{
}

\Verse{45}
{
    \zA{}
}
{
    \eA{Hsiang Ling drew near and passed it under inspection. "No wonder," she
        exclaimed, "it looks so handsome! It's verily woven with peacock's
        feathers." "What about peacock's feathers?" Hsiang-yün laughed. "It's made
        of the feathers plucked from the heads of wild ducks. This is a clear sign
        that our worthy ancestor is fond of you, for with all her love for Pao-yü,
        she hasn't given it to him to wear." "Truly does the proverb say: 'that
        every human being has his respective lot.'" Pao-ch'ai smiled. "Nothing ever
        was further from my thoughts than that she would, at this juncture, drop on
        the scene! Come she may, but here she also gets our dear ancestor to lavish
        such love on her!" "Unless you stay with our worthy senior," Hsiang-yün
        said, "do come into the garden.}
}
{
}

\Verse{46}
{
    \zA{}
}
{
    \eA{You may romp and laugh and eat and drink as much as you like in these two
        places. But when you get over to Madame Hsing's rooms, talk and joke with
        her, if she be at home, to your heart's content; it won't matter if you
        tarry ever so long. But should she not be in, don't put your foot inside;
        for the inmates are many in those rooms and their hearts are evil. All
        they're up to is to do us harm." These words much amused Pao-ch'ai,
        Pao-ch'in, Hsiang-Ling, Ying Erh and the others present. "Were one to say,"
        Pao-ch'ai smiled, "that you're heartless, (it wouldn't do); for you've got a
        heart. But despite your having a heart, your tongue is, in fact, a little
        too outspoken! You should really to-day acknowledge this Ch'in Erh of ours
        as your own sister!"}
}
{
}

\Verse{47}
{
    \zA{}
}
{
    \eA{"This article of clothing," Hsiang-yün laughed, casting another glance at
        Pao-ch'in, "is only meet for her to wear. It wouldn't verily look well on
        any one else." Saying this, she espied Hu Po enter the room. "Our old
        mistress," she put in smiling, "bade me tell you, Miss Pao-ch'ai, not to
        keep too strict a check over Miss Ch'in, for she's yet young; that you
        should let her do as she pleases, and that whatever she wants you should ask
        for, and not be afraid." Pao-ch'ai hastily jumped to her feet and signified
        her obedience. Pushing Pao-ch'in, she laughed. "Even you couldn't tell
        whence this piece of good fortune hails from," she said. "Be off now; for
        mind, we might hurt your feelings. I can never believe myself so inferior to
        you!"}
}
{
}

\Verse{48}
{
    \zA{}
}
{
    \eA{As she spoke, Pao-yü and Tai-yü walked in. But as Pao-ch'ai continued to
        indulge in raillery to herself, "Cousin Pao," Hsiang-yün smilingly
        remonstrated, "you may, it's true, be jesting, but what if there were any
        one to entertain such ideas in real earnest?" "If any one took things in
        earnest," Hu Po interposed laughing, "why, she'd give offence to no one else
        but to him." Pointing, as she uttered this remark, at Pao-yü. "He's not that
        sort of person!" Pao-ch'ai and Hsiang-yün simultaneously ventured, with a
        significant smile. "If it isn't he," Hu Po proceeded still laughing, "it's
        she." Turning again her finger towards Tai-yü. Hsiang-yün expressed not a
        word by way of rejoinder. "That's still less likely," Pao-ch'ai smiled, "for
        my cousin is like her own sister;}
}
{
}

\Verse{49}
{
    \zA{}
}
{
    \eA{and she's far fonder of her than of me. How could she therefore take
        offence? Do you credit that nonsensical trash uttered by Yün-erh! Why what
        good ever comes out of that mouth of hers?" Pao-yü was ever well aware that
        Tai-yü was gifted with a somewhat mean disposition. He had not however as
        yet come to learn anything of what had recently transpired between Tai-yü
        and Pao-ch'ai. He was therefore just giving way to fears lest his
        grandmother's fondness for Pao-ch'in should be the cause of her feeling
        dejected.}
}
{
}

\Verse{50}
{
    \zA{}
}
{
    \eA{But when he now heard the remarks passed by Hsiang-yün, and the rejoinders
        made, on the other hand, by Pao-ch'ai, and, when he noticed how different
        Tai-yü's voice and manner were from former occasions, and how they actually
        bore out Pao-ch'ai's insinuation, he was at a great loss how to solve the
        mystery. "These two," he consequently pondered, "were never like this
        before! From all I can now see, they're, really, a hundred times far more
        friendly than any others are!" But presently he also observed Lin Tai-yü
        rush after Pao-ch'in, and call out 'Sister,' and, without even making any
        allusion to her name or any mention to her surname, treat her in every
        respect, just as if she were her own sister. This Pao-ch'in was young and
        warm-hearted.}
}
{
}

\Verse{51}
{
    \zA{}
}
{
    \eA{She was naturally besides of an intelligent disposition. She had, from her
        very youth up, learnt how to read and how to write. After a stay, on the
        present occasion, of a couple of days in the Chia mansion, she became
        acquainted with nearly every inmate. And as she saw that the whole bevy of
        young ladies were not of a haughty nature, and that they kept on friendly
        terms with her own cousin, she did not feel disposed to treat them with any
        discourtesy. But she had likewise found out for herself that Lin Tai-yü was
        the best among the whole lot, so she started with Tai-yü, more than with any
        one else, a friendship of unusual fervour. This did not escape Pao-yü's
        notice;}
}
{
}

\Verse{52}
{
    \zA{}
}
{
    \eA{but all he could do was to secretly give way to amazement. Shortly, however,
        Pao-ch'ai and her cousin repaired to Mrs. Hsüeh's quarters. Hsiang-yün then
        betook herself to dowager lady Chia's apartments, while Lin Tai-yü returned
        to her room and lay down to rest. Pao-yü thereupon came to look up Tai-yü.
        "Albeit I've read the 'Record of the Western Side-room,'" he smiled, "and
        understood a few passages of it, yet when I quoted some in order to make you
        laugh, you flew into a huff! But I now remember that there is, indeed, a
        passage, which is not intelligible to me; so let me quote it for you to
        explain it for me!"}
}
{
}

\Verse{53}
{
    \zA{}
}
{
    \eA{Hearing this, Tai-yü immediately concluded that his words harboured some
        secret meaning, so putting on a smile, "Recite it and let me hear it," she
        said. "In the 'Confusion' chapter," Pao-yü laughingly began, "there's a line
        couched in most beautiful language. It's this: 'What time did Meng Kuang
        receive Liang Hung's candlestick?' (When did you and Pao-ch'ai get to be
        such friends?) These five characters simply bear on a stock story; but to
        the credit of the writer be it, the question contained in the three empty
        words representing, 'What time' is set so charmingly! When did she receive
        it? Do tell me!" At this inquiry, Tai-yü too could not help laughing. "The
        question was originally nicely put," she felt urged to rejoin with a laugh.}
}
{
}

\Verse{54}
{
    \zA{}
}
{
    \eA{"But though the writer sets it gracefully, you ask it likewise with equal
        grace!" "At one time," Pao-yü. observed, "all you knew was to suspect that I
        (was in love with Pao-ch'ai); and have you now no faults to find?" "Who ever
        could have imagined her such a really nice girl!" Tai-yü smiled. "I've all
        along thought her full of guile!" And seizing the occasion, she told Pao-yü
        with full particulars how she had, in the game of forfeits, made an improper
        quotation, and what advice Pao-ch'ai had given her on the subject; how she
        had even sent her some birds' nests, and what they had said in the course of
        the chat they had had during her illness. Pao-yü then at length came to see
        why it was that such a warm friendship had sprung up between them.}
}
{
}

\Verse{55}
{
    \zA{}
}
{
    \eA{"To tell you the truth," he consequently remarked smilingly, "I was just
        wondering when Meng Kuang had received Liang Hung's candlestick; and, lo,
        you, indeed, got it, when a mere child and through some reckless talk, (and
        your friendship was sealed)." As the conversation again turned on Pao-ch'in,
        Tai-yü recalled to mind that she had no sister, and she could not help
        melting once more into tears. Pao-yü hastened to reason with her. "This is
        again bringing trouble upon yourself!" he argued. "Just see how much thinner
        you are this year than you were last; and don't you yet look after your
        health? You deliberately worry yourself every day of your life. And when
        you've had a good cry, you feel at last that you've acquitted yourself of
        the duties of the day."}
}
{
}

\Verse{56}
{
    \zA{}
}
{
    \eA{"Of late," Tai-yü observed, drying her tears, "I feel sore at heart. But my
        tears are scantier by far than they were in years gone by. With all the
        grief and anguish, which gnaw my heart, my tears won't fall plentifully."
        "This is because weeping has become a habit with you," Pao-yü added. "But
        though you fancy to yourself that it is so, how can your tears have become
        scantier than they were?" While arguing with her, he perceived a young
        waiting-maid, attached to his room, bring him a red felt wrapper. "Our
        senior mistress, lady Chia Chu," she went on, "has just sent a servant to
        say that, as it snows, arrangements should be made for inviting people
        to-morrow to write verses." But hardly was this message delivered, than they
        saw Li Wan's maid enter, and invite Tai-yü to go over.}
}
{
}

\Verse{57}
{
    \zA{}
}
{
    \eA{Pao-yü then proposed to Tai-yü to accompany him, and together they came to
        the Tao Hsiang village. Tai-yü changed her shoes for a pair of low shoes
        made of red scented sheep skin, ornamented with gold, and hollowed clouds.
        She put on a deep red crape cloak, lined with white fox fur; girdled herself
        with a lapis-lazuli coloured sash, decorated with bright green double rings
        and four sceptres; and covered her head with a hat suitable for rainy
        weather. After which, the two cousins trudged in the snow, and repaired to
        this side of the mansion. Here they discovered the young ladies assembled,
        dressed all alike in deep red felt or camlet capes, with the exception of Li
        Wan, who was clad in a woollen jacket, buttoning in the middle.}
}
{
}

\Verse{58}
{
    \zA{}
}
{
    \eA{Hsüeh Pao-ch'ai wore a pinkish-purple twilled pelisse, lined with foreign
        'pa' fur, worked with threads from abroad, and ornamented with double
        embroidery. Hsing Chou-yen was still attired in an old costume, she
        ordinarily used at home, without any garment for protection against the
        rain. Shortly, Shih Hsiang-yün arrived. She wore the long pelisse, given her
        by dowager lady Chia, which gave warmth both from the inside and outside, as
        the top consisted of martin-head fur, and the lining of the long-haired coat
        of the dark grey squirrel. On her head, she had a deep red woollen hood,
        made \_á la\_ Chao Chün, with designs of clouds scooped out on it. This was
        lined with gosling-yellow, gold-streaked silk.}
}
{
}

\Verse{59}
{
    \zA{}
}
{
    \eA{Round her neck, she had a collar of sable fur. "Just see here!" Tai-yü was
        the first to shout with a laugh. "Here comes Sun Hsing-che the
        'monkey-walker!' Lo, like him, she holds a snow cloak, and purposely puts on
        the air of a young bewitching ape!" "Look here, all of you!" Hsiang-yün
        laughed. "See what I wear inside!" So saying, she threw off her cloak. This
        enabled them to notice that she wore underneath a half-new garment with
        three different coloured borders on the collar and cuffs, consisting of a
        short pelisse of russet material lined with ermine and ornamented with
        dragons embroidered in variegated silks whose coils were worked with golden
        threads. The lapel was narrow. The sleeves were short. The folds buttoned on
        the side.}
}
{
}

\Verse{60}
{
    \zA{}
}
{
    \eA{Under this, she had a very short light-red brocaded satin bodkin, lined with
        fur from foxes' ribs. Round her waist was lightly attached a many-hued
        palace sash, with butterfly knots and long tassels. On her feet, she too
        wore a pair of low shoes made of deer leather. Her waist looked more than
        ever like that of a wasp, her back like that of the gibbon. Her bearing
        resembled that of a crane, her figure that of a mantis. "Her weak point,"
        they laughed unanimously, "is to get herself up to look like a young masher.
        But she does, there's no denying, cut a much handsomer figure like this,
        than when she's dressed up like a girl!" "Lose no time," Hsiang-yün smiled,
        "in deliberating about writing verses, for I'd like to hear who is to stand
        treat."}
}
{
}

\Verse{61}
{
    \zA{}
}
{
    \eA{"According to my idea," Li Wan chimed in, "I think that as the legitimate
        day, which was yesterday, has gone by, it would be too long to wait for
        another proper date. As luck would have it, it's snowing again to-day, so
        won't it be well to raise contributions among ourselves and have a meeting?
        We'll thus be able to give the visitors a greeting; and to get an
        opportunity of writing a few verses. But what are your views on the
        subject?" "This proposal is excellent!" Pao-yü was the first to exclaim.
        "The only thing is that it's too late to-day; and if it clears up by
        to-morrow, there will be really no fun." "It isn't likely," cried out the
        party with one voice, "that this snowy weather will clear up. But even
        supposing it does, the snow which will fall during this night will be
        sufficient for our enjoyment."}
}
{
}

\Verse{62}
{
    \zA{}
}
{
    \eA{"This place of mine is nice enough, it's true," Li Wan added, "yet it isn't
        up to the Lu Hsüeh Pavilion. I've already therefore despatched workmen to
        raise earthen couches, so that we should all be able to sit round the fire
        and compose our verses. Our venerable senior, I fancy, is not sure about
        caring to join us. Besides, this is only a small amusement between ourselves
        so if we just let that hussy Feng know something about it, it will be quite
        enough. A tael from each of you will be ample, but send your money to me
        here! As regards Hsiang Ling, Pao-ch'in, Li Wen, Li Ch'i and Chou-yen, the
        five of them, we needn't count them. Neither need we include the two girls
        of our number, who are ill;}
}
{
}

\Verse{63}
{
    \zA{}
}
{
    \eA{nor take into account the four girls who've asked for leave. If you will let
        me have your four shares, I'll undertake to see that five or six taels be
        made to suffice." Pao-ch'ai and the others without exception signified their
        acquiescence. They consequently proceeded to propose the themes and to fix
        upon the rhymes. "I've long ago," smiled Li Wan, "settled them in my own
        mind, so tomorrow at the proper time you'll really know all about them." At
        the conclusion of this remark, they indulged in another chat on irrelevant
        topics; and this over, they came into old lady Chia's quarters. Nothing of
        any note transpired during the course of that day.}
}
{
}

\Verse{64}
{
    \zA{}
}
{
    \eA{At an early hour on the morrow, Pao-yü—for he had been looking forward with
        such keen expectation to the coming event that he had found it impossible to
        have any sleep during the night,—jumped out of bed with the first blush of
        dawn. Upon raising his curtain and looking out, he observed that, albeit the
        doors and windows were as yet closed, a bright light shone on the lattice
        sufficient to dazzle the eyes, and his mind began at once to entertain
        misgivings, and to feel regrets, in the assurance that the weather had
        turned out fine, and that the sun had already risen.}
}
{
}

\Verse{65}
{
    \zA{}
}
{
    \eA{In a hurry, he simultaneously sprung to his feet, and flung the window-frame
        open, then casting a glance outside, from within the glass casement, he
        realised that it was not the reflection of the sun, but that of the snow,
        which had fallen throughout the night to the depth of over a foot, and that
        the heavens were still covered as if with twisted cotton and unravelled
        floss. Pao-yü got, by this time, into an unusual state of exhilaration.
        Hastily calling up the servants, and completing his ablutions, he robed
        himself in an egg-plant-coloured camlet, fox-fur lined pelisse; donned a
        short-sleeved falconry surtout ornamented with water dragons; tied a sash
        round his waist; threw over his shoulders a fine bamboo waterproof; covered
        his head with a golden rattan rain-hat;}
}
{
}

\Verse{66}
{
    \zA{}
}
{
    \eA{put on a pair of 'sha t'ang' wood clogs, and rushed out with precipitate
        step towards the direction of the Lu Hsüeh Pavilion. As soon as he sallied
        out of the gate of the courtyard, he gazed on all four quarters. No trace
        whatever of any other colour (but white) struck his eye. In the distance
        stood the green fir-trees and the kingfisherlike bamboos. They too looked,
        however, as if they were placed in a glass bowl. Forthwith he wended his way
        down the slope and trudged along the foot of the hill. But the moment he
        turned the bend, he felt a whiff of cold fragrance come wafted into his
        nostrils. Turning his head, he espied ten and more red plum trees, over at
        Miao Yü's in the Lung Ts'ui monastery. They were red like very rouge.}
}
{
}

\Verse{67}
{
    \zA{}
}
{
    \eA{And, reflecting the white colour of the snow, they showed off their beauty
        to such an extraordinary degree as to present a most pleasing sight. Pao-yü
        quickly stood still, and gazed, with all intentness, at the landscape for a
        time. But just as he was proceeding on his way, he caught sight of some one
        on the "Wasp waist" wooden bridge, advancing in his direction, with an
        umbrella in hand. It was the servant, despatched by Li Wan, to request lady
        Peng to go over. On his arrival in the Lu Hsüeh pavilion, Pao-yü found the
        maids and matrons engaged in sweeping away the snow and opening a passage.
        This Lu Hsüeh (Water-rush snow) pavilion was, we might explain, situated on
        a side hill, in the vicinity of a stream and spanned the rapids formed by
        it.}
}
{
}

\Verse{68}
{
    \zA{}
}
{
    \eA{The whole place consisted of several thatched roofs, mud walls, side fences,
        bamboo lattice windows and pushing windows, out of which fishing-lines could
        be conveniently dropped. On all four sides flourished one mass of reeds,
        which concealed the single path out of the pavilion. Turning and twisting,
        he penetrated on his way through the growth of reeds until he reached the
        spot where stretched the bamboo bridge leading to the Lotus Fragrance
        Arbour.}
}
{
}

\Verse{69}
{
    \zA{}
}
{
    \eA{The moment the maids and matrons saw him approach with his
        waterproof-wrapper thrown over his person and his rain-hat on his head, they
        with one voice laughed, "We were just remarking that what was lacking was a
        fisherman, and lo, now we've got everything that was wanted! The young
        ladies are coming after their breakfast; you're in too impatient a mood!" At
        these words, Pao-yü had no help but to retrace his footsteps. As soon as he
        reached the Hsin Tang pavilion, he perceived T'an Ch'un, issuing from the
        Ch'iu Shuang Study, wrapped in a deep red woollen waterproof, and a 'Kuan
        Yin' hood on her head, supporting herself on the arm of a young maid. Behind
        her, followed a married woman, holding a glazed umbrella made of green
        satin.}
}
{
}

\Verse{70}
{
    \zA{}
}
{
    \eA{Pao-yü knew very well that she was on her way to his grandmother's, so
        speedily halting by the side of the pavilion, he waited for her to come up.
        The two cousins then left the garden together, and betook themselves to the
        front part of the mansion. Pao-ch'in was at the time in the inner
        apartments, combing her hair, washing her hands and face and changing her
        apparel. Shortly, the whole number of girls arrived. "I feel peckish!"
        Pao-yü shouted; and again and again he tried to hurry the meal. It was with
        great impatience that he waited until the eatables could be laid on the
        table. One of the dishes consisted of kid, boiled in cow's milk. "This is
        medicine for us, who are advanced in years," old lady Chia observed.
        "They're things that haven't seen the light! The pity is that you young
        people can't have any.}
}
{
}

\Verse{71}
{
    \zA{}
}
{
    \eA{There's some fresh venison to-day as an extra course, so you'd better wait
        and eat some of that!" One and all expressed their readiness to wait. Pao-yü
        however could not delay having something to eat. Seizing a cup of tea, he
        soaked a bowlful of rice, to which he added some meat from a pheasant's leg,
        and gobbled it down in a scramble. "I'm well aware," dowager lady Chia said,
        "that as you're up to something again to-day, you people have no mind even
        for your meal. Let them keep," she therefore cried, "that venison for their
        evening repast!" "What an idea!" lady Feng promptly put in. "We'll have
        enough with what remains of it." Shih Hsiang-yün thereupon consulted with
        Pao-yü.}
}
{
}

\Verse{72}
{
    \zA{}
}
{
    \eA{"As there's fresh venison," she said, "wouldn't it be nice to ask for a
        haunch and take it into the garden and prepare it ourselves? We'll thus be
        able to sate our hunger, and have some fun as well." At this proposal,
        Pao-yü actually asked lady Feng to let them have a haunch, and he bade a
        matron carry it into the garden. Presently, they all got up from table.
        After a time, they entered the garden and came in a body to the Lu Hsüeh
        pavilion to hear Li Wan give out the themes, and fix upon the rhymes. But
        Hsiang-yün and Pao-yü were the only two of whom nothing was seen. "Those
        two," Tai-yü observed, "can't get together! The moment they meet, how much
        trouble doesn't arise! They must surely have now gone to hatch their plans
        over that haunch of venison."}
}
{
}

\Verse{73}
{
    \zA{}
}
{
    \eA{These words were still on her lips when she saw 'sister-in-law' Li coming
        also to see what the noise was all about. "How is it," she then inquired of
        Li Wan, "that that young fellow, with the jade, and that girl, with the
        golden unicorn round her neck, both of whom are so cleanly and tidy, and
        have besides ample to eat, are over there conferring about eating raw meat?
        There they are chatting, saying this and saying that; but I can't see how
        meat can be eaten raw!" This remark much amused the party. "How dreadful!"
        they exclaimed, "Be quick and bring them both here!" "All this fuss," Tay-yü
        smiled, "is the work of that girl Yün. I'm not far off again in my
        surmises." Li Wan went out with precipitate step in search of the cousins.}
}
{
}

\Verse{74}
{
    \zA{}
}
{
    \eA{"If you two are bent upon eating raw meat," she cried, "I'll send you over
        to our old senior's; you can do so there. What will I care then if you have
        a whole deer raw and make yourselves ill over it? It won't be any business
        of mine. But it's snowing hard and it's bitterly cold, so be quick and go
        and write some verses for me and be off!" "We're doing nothing of the kind,"
        Pao-yü hastily rejoined. "We're going to eat some roasted meat." "Well, that
        won't matter!" Li Wan observed. And seeing the old matrons bring an iron
        stove, prongs and a gridiron of iron wire, "Mind you don't cut your hands,"
        Li Wan resumed, "for we won't have any crying!" This remark concluded, she
        walked in.}
}
{
}

\Verse{75}
{
    \zA{}
}
{
    \eA{Lady Feng had sent P'ing Erh from her quarters to announce that she was
        unable to come, as the issue of the customary annual money gave her just at
        present, plenty to keep her busy. Hsiang-yün caught sight of P'ing Erh and
        would not let her go on her errand. But P'ing Erh too was fond of amusement,
        and had ever followed lady Feng everywhere she went, so, when she perceived
        what fun was to be got, and how merrily they joked and laughed, she felt
        impelled to take off her bracelets (and to join them). The trio then pressed
        round the fire; and P'ing Erh wanted to be the first to roast three pieces
        of venison to regale themselves with.}
}
{
}

\Verse{76}
{
    \zA{}
}
{
    \eA{On the other side, Pao-ch'ai and Tai-yü had, even in ordinary times, seen
        enough of occasions like the present. They did not therefore think it
        anything out of the way; but Pao-ch'in and the other visitors, inclusive of
        'sister-in-law' Li, were filled with intense wonder. T'an Ch'un had, with
        the help of Li Wan, and her companions, succeeded by this time in choosing
        the subjects and rhymes. "Just smell that sweet fragrance," T'an Ch'un
        remarked. "One can smell it even here! I'm also going to taste some." So
        speaking, she too went to look them up. But Li Wan likewise followed her
        out. "The guests are all assembled," she observed. "Haven't you people had
        enough as yet?" While Hsiang-yün munched what she had in her month, she
        replied to her question.}
}
{
}

\Verse{77}
{
    \zA{}
}
{
    \eA{"Whenever," she said, "I eat this sort of thing, I feel a craving for wine.
        It's only after I've had some that I shall be able to rhyme. Were it not for
        this venison, I would to-day have positively been quite unfit for any
        poetry." As she spoke, she discerned Pao-ch'in, standing and laughing
        opposite to her, in her duck-down garment. "You idiot," Hsiang-yün
        laughingly cried, "come and have a mouthful to taste." "It's too filthy!"
        Pao-ch'in replied smiling. "You go and try it." Pao-ch'ai added with a
        laugh. "It's capital! Your cousin Lin is so very weak that she couldn't
        digest it, if she had any. Otherwise she too is very fond of this." Upon
        hearing this, Pao-ch'in readily crossed over and put a piece in her mouth;
        and so good did she find it that she likewise started eating some of it.}
}
{
}

\Verse{78}
{
    \zA{}
}
{
    \eA{In a little time, however, lady Feng sent a young maid to call P'ing Erh.
        "Miss Shih," P'ing Erh explained, "won't let me go. So just return ahead of
        me." The maid thereupon took her leave; but shortly after they saw lady Feng
        arrive; she too with a wrapper over her shoulders. "You're having," she
        smiled; "such dainties to eat, and don't you tell me?" Saying this, she also
        drew near and began to eat. "Where has this crowd of beggars turned up
        from?" Tai-yü put in with a laugh. "But never mind, never mind! Here's the
        Lu Hsüeh pavilion come in for this calamity to-day, and, as it happens, it's
        that chit Yün by whom it has been polluted! But I'll have a good cry for the
        Lu Hsüeh pavilion." Hsiang-yün gave an ironical smile. "What do you know?"
        she exclaimed. "A genuine man of letters is naturally refined.}
}
{
}

\Verse{79}
{
    \zA{}
}
{
    \eA{But as for the whole lot of you, your poor and lofty notions are all a sham!
        You are most loathsome! We may now be frowzy and smelly, as we munch away
        lustily with our voracious appetites, but by and bye we'll prove as refined
        as scholars, as if we had cultured minds and polished tongues." "If by and
        bye," Pao-ch'ai laughingly interposed, "the verses you compose are not worth
        anything, I'll tug out that meat you've eaten, and take some of these
        snow-buried weeds and stuff you up with. I'll thus put an end to this evil
        fortune!" While bandying words, they finished eating. For a time, they
        busied themselves with washing their hands. But when P'ing Erh came to put
        on her bracelets, she found one missing.}
}
{
}

\Verse{80}
{
    \zA{}
}
{
    \eA{She looked in a confused manner, at one time to the left, at another to the
        right; now in front of her, and then behind her for ever so long, but not a
        single vestige of it was visible. One and all were therefore filled with
        utter astonishment. "I know where this bracelet has gone to;" lady Feng
        suggested smilingly. "But just you all go and attend to your poetry. We too
        can well dispense with searching for it, and repair to the front. Before
        three days are out, I'll wager that it turns up. What verses are you writing
        to-day?" continuing she went on to inquire. "Our worthy senior says that the
        end of the year is again nigh at hand, and that in the first moon some more
        conundrums will have to be devised to be affixed on lanterns, for the
        recreation of the whole family."}
}
{
}

\Verse{81}
{
    \zA{}
}
{
    \eA{"Of course we'll have to write a few," they laughingly rejoined, upon
        hearing her remarks. "We forgot all about it. Let's hurry up now, and
        compose a few fine ones, so as to have them ready to enjoy some good fun in
        the first moon." Speaking the while, they came in a body into the room with
        the earthen couches, where they found the cups, dishes and eatables already
        laid out in readiness. On the walls had been put up the themes, metre, and
        specimen verses. Pao-yü and Hsiang-yün hastened to examine what was written.
        They saw that they had to take for a theme something on the present scenery
        and indite a stanza with antithetical pentameter lines; that the word
        'hsiao,' second (in the book of metre), had been fixed upon as a rhyme;}
}
{
}

\Verse{82}
{
    \zA{}
}
{
    \eA{but that there was, below that, no mention, as yet, made of any precedence.
        "I can't write verses very well," Li Wan pleaded, "so all I'll do will be to
        devise three lines, and the one, who'll finish the task first, we'll have
        afterwards to pair them." "We should, after all," Pao-ch'ai urged, "make
        some distinction with regard to order." But, reader, if you entertain any
        desire to know the sequel, peruse the particulars recorded in the chapter
        that follows.}
}
{
}
